\scanpage{24}\mspage{22}
Je dis que cette composition est associative. Cela
résulte du fait que si on \unsure{a} trois chemins (d'indices
\uline{pairs} !)
\diagram{\begin{tikzcd}[ampersand replacement=\&]
x \arrow[r,no head,"c"] \& y \arrow[r,no head,"c'"] \& \deleted{\unsure{t}}z \arrow[r,"c''"] \& \deleted{\unsure{z}}t
\end{tikzcd}},
on a \unsure{$^{*}$} (sans avoir à effectuer des \unsure{homotopies}\ldots)
associativité stricte
\[
c''\mathbin{.}(c'\mathbin{.}c)=(c''\mathbin{.}c')\mathbin{.}c\quad .
\]
Je dis qu'il y a des identités. On définit
\[
1_{\unsure{x|z}}\in\operatorname{Hom}_{\unsure{\Sigma|\overline{Z}|\overline{\Sigma}|Z}}(x,x)
\]
par le chemin \deleted{de longueur} \inserted{d'indice} $0$ (qui est pair !)
i.e. par le diagramme réduit à $x$, sans flèche
(ou si on veut, en interprétant le diagramme comme un
foncteur $e\to X$, il y a la flèche identique
sous-entendue). Nos raisonnements \unsure{n'excluaient|n'excluent}
pas le cas $n=0$ ou $m=0$, et il est \deleted{évident} \inserted{tautologique} que
pour un chemin
\diagram{\begin{tikzcd}[ampersand replacement=\&]x\arrow[r,no head,"c"]\&y\end{tikzcd}}
on a
\[
c\mathbin{.}1_x=1_y\mathbin{.}c=c\quad .
\]
Donc on a bien défini une catégorie, que je
vais appeler $X^{\deleted{\unsure{*}}}_{\Sigma}$. Je vais définir un
foncteur canonique
\[
\deleted{\unsure{Q|\Theta}}\unsure{\varphi|\Theta|Q}\colon X\longrightarrow X_{\Sigma}
\]
qui sera l'identité sur les objets, il faut le
définir sur les flèches
\[
\operatorname{Hom}_X(x,y)\unsure{\mathord{-}|\to}\operatorname{Hom}_{X_{\Sigma}}(x,y)=\pi_0\mathrm{Ch}_{\Sigma}(x,y)
\]
Malheureusement, une flèche
\diagram{\begin{tikzcd}[ampersand replacement=\&]x\arrow[r,no head,"u"]\&y\end{tikzcd}}
% archive sheet number: 23
